\documentclass[a4paper,twoside]{article}

\usepackage{morewrites}

\usepackage[svgnames,dvipsnames,x11names]{xcolor}
\usepackage{graphicx}
\usepackage[english]{babel}
\usepackage[colorlinks=true, linkcolor=NavyBlue, urlcolor=NavyBlue, citecolor=NavyBlue]{hyperref}
\usepackage[a4paper]{geometry}
\usepackage{ts-fonts}
\usepackage{float}
\usepackage{graphicx}
\usepackage{seqsplit}

\usepackage{lastpage}
\usepackage{refcount}
\makeatletter
\def\ps@plain{
\let\@oddhead\@empty
\let\@evenhead\@empty
\def\@oddfoot{
\hfil
\ifnum\getpagerefnumber{LastPage}>1\relax
\thepage
\fi
\hfil}
\let\@evenfoot\@oddfoot
}
\makeatother

\title{Author}
\author{}\date{}
\begin{document}
\maketitle
\begin{figure}[H]
\centering
\includegraphics[width=0.5\linewidth]{52489316.jpg}
\caption{Yves Chevallier}
\end{figure}
I'm Yves Chevallier, the creator and maintainer of TeXSmith. I work full-time as an associate professor at \href{https://www.heig-vd.ch}{HEIG-VD} in Switzerland, where I teach software engineering. My research interests revolve around software architecture, meta-generation, and---naturally---documentation tooling. Before drifting into academia, I trained as an electronics engineer with a focus on real-time embedded systems for motion-control applications.

I originally built TeXSmith mostly for fun---and partly to see how far AI-assisted tooling can be pushed. Judging by the pace of progress, this was probably a terrible idea: my job may well be obsolete in a few years.
\end{document}
